% Bài: L12C1 Bài 2. Nội năng. Định luật I của nhiệt động lực học · Tính độ biến thiên nội năng trong quá trình nhiệt động lực học
% ID: L12C1B2-01-DS
% Mức: NB
% ID: L12C1B2-01-DS


% ID: L12C1B2-16-DS
% Mức: NB
% ID: L12C1B2-74-TL
% ID: L12C1B2-74-TL
% Mức: NB
% ID: L12C1B3-100-TL
% ID: L12C1B3-100-TL
% Mức: NB
% ID: L12C1B2-81-TL
%=== Câu 41 ===%
\dangbt{Tính độ biến thiên nội năng trong quá trình nhiệt động lực học}
\begin{ex}
% ID: L12C1B2-81-TL
% Mức: NB
Xét các phát biểu sau về dạng “Tính độ biến thiên nội năng trong quá trình nhiệt động lực học”
{\True Nếu hệ nhận công 200 $J$ và tỏa 50 $J$ thì ΔU= $150\,\text{J}$.}
{Nếu A=100 $J$ và Q=-40 $J$ thì ΔU=- $140\,\text{J}$.}
{\True Khi giải dạng “Tính độ biến thiên nội năng trong quá trình nhiệt động lực học”, cần xác định đúng trạng thái đầu, trạng thái cuối và quy ước dấu hoặc điều kiện áp dụng.}
{Có thể bỏ qua mọi dữ kiện về khối lượng, nhiệt độ và hiệu suất trong mọi bài toán thuộc dạng này.}
\loigiai{
\begin{itemchoice}
\item (Đúng) Nếu hệ nhận công 200 $J$ và tỏa 50 $J$ thì ΔU= $150\,\text{J}$. (Vì): Nếu hệ nhận công 200 $J$ và tỏa 50 $J$ thì ΔU= $150\,\text{J}$. b) S: Nếu A=100 $J$ và Q=-40 $J$ thì ΔU=- $140\,\text{J}$. c) Đ: Khi giải dạng “Tính độ biến thiên nội năng trong quá trình nhiệt động lực học”, cần xác định đúng trạng thái đầu, trạng thái cuối và quy ước dấu hoặc điều kiện áp dụng. d) S: Có thể bỏ qua mọi dữ kiện về khối lượng, nhiệt độ và hiệu suất trong mọi bài toán thuộc dạng này. Kết luận dựa trên định nghĩa, điều kiện áp dụng và quy ước trong SGK KNTT
\item (Sai) Nếu A=100 $J$ và Q=-40 $J$ thì ΔU=- $140\,\text{J}$.
\item (Đúng) Khi giải dạng “Tính độ biến thiên nội năng trong quá trình nhiệt động lực học”, cần xác định đúng trạng thái đầu, trạng thái cuối và quy ước dấu hoặc điều kiện áp dụng.
\item (Sai) Có thể bỏ qua mọi dữ kiện về khối lượng, nhiệt độ và hiệu suất trong mọi bài toán thuộc dạng này.
\end{itemchoice}
}
\end{ex}

% ID: L12C1B2-17-DS
% Mức: TH
% ID: L12C1B2-17-DS
% ID: L12C1B2-17-DS
% Mức: TH
% ID: L12C1B3-101-DS
% ID: L12C1B3-101-DS
% Mức: TH
% ID: L12C1B2-83-DS
%=== Câu 43 ===%
\begin{ex}
% ID: L12C1B2-83-DS
% Mức: TH
Xét các phát biểu sau về dạng “Tính độ biến thiên nội năng trong quá trình nhiệt động lực học”
\choiceTF
{Nếu A=100 $J$ và Q=-40 $J$ thì ΔU=- $140\,\text{J}$.}
{\True Khi giải dạng “Tính độ biến thiên nội năng trong quá trình nhiệt động lực học”, cần xác định đúng trạng thái đầu, trạng thái cuối và quy ước dấu hoặc điều kiện áp dụng.}
{Có thể bỏ qua mọi dữ kiện về khối lượng, nhiệt độ và hiệu suất trong mọi bài toán thuộc dạng này.}
{\True Nếu hệ nhận công 200 $J$ và tỏa 50 $J$ thì ΔU= $150\,\text{J}$.}
\loigiai{
\begin{itemchoice}
\item (Sai) Nếu A=100 $J$ và Q=-40 $J$ thì ΔU=- $140\,\text{J}$. (Vì): Nếu A=100 $J$ và Q=-40 $J$ thì ΔU=- $140\,\text{J}$. b) Đ: Khi giải dạng “Tính độ biến thiên nội năng trong quá trình nhiệt động lực học”, cần xác định đúng trạng thái đầu, trạng thái cuối và quy ước dấu hoặc điều kiện áp dụng. c) S: Có thể bỏ qua mọi dữ kiện về khối lượng, nhiệt độ và hiệu suất trong mọi bài toán thuộc dạng này. d) Đ: Nếu hệ nhận công 200 $J$ và tỏa 50 $J$ thì ΔU= $150\,\text{J}$. Kết luận dựa trên định nghĩa, điều kiện áp dụng và quy ước trong SGK KNTT
\item (Đúng) Khi giải dạng “Tính độ biến thiên nội năng trong quá trình nhiệt động lực học”, cần xác định đúng trạng thái đầu, trạng thái cuối và quy ước dấu hoặc điều kiện áp dụng.
\item (Sai) Có thể bỏ qua mọi dữ kiện về khối lượng, nhiệt độ và hiệu suất trong mọi bài toán thuộc dạng này.
\item (Đúng) Nếu hệ nhận công 200 $J$ và tỏa 50 $J$ thì ΔU= $150\,\text{J}$.
\end{itemchoice}
}
\end{ex}

% ID: L12C1B2-18-DS
% Mức: TH
% ID: L12C1B2-18-DS
% ID: L12C1B2-18-DS
% Mức: TH
% ID: L12C1B3-102-DS
% ID: L12C1B3-102-DS
% Mức: TH
% ID: L12C1B2-84-DS
%=== Câu 45 ===%
\begin{ex}
% ID: L12C1B2-84-DS
% Mức: TH
Xét các phát biểu sau về dạng “Tính độ biến thiên nội năng trong quá trình nhiệt động lực học”
\choiceTF
{Nếu A=100 $J$ và Q=-40 $J$ thì ΔU=- $140\,\text{J}$.}
{\True Khi giải dạng “Tính độ biến thiên nội năng trong quá trình nhiệt động lực học”, cần xác định đúng trạng thái đầu, trạng thái cuối và quy ước dấu hoặc điều kiện áp dụng.}
{Có thể bỏ qua mọi dữ kiện về khối lượng, nhiệt độ và hiệu suất trong mọi bài toán thuộc dạng này.}
{\True Nếu hệ nhận công 200 $J$ và tỏa 50 $J$ thì ΔU= $150\,\text{J}$.}
\loigiai{
\begin{itemchoice}
\item (Sai) Nếu A=100 $J$ và Q=-40 $J$ thì ΔU=- $140\,\text{J}$. (Vì): Nếu A=100 $J$ và Q=-40 $J$ thì ΔU=- $140\,\text{J}$. b) Đ: Khi giải dạng “Tính độ biến thiên nội năng trong quá trình nhiệt động lực học”, cần xác định đúng trạng thái đầu, trạng thái cuối và quy ước dấu hoặc điều kiện áp dụng. c) S: Có thể bỏ qua mọi dữ kiện về khối lượng, nhiệt độ và hiệu suất trong mọi bài toán thuộc dạng này. d) Đ: Nếu hệ nhận công 200 $J$ và tỏa 50 $J$ thì ΔU= $150\,\text{J}$. Kết luận dựa trên định nghĩa, điều kiện áp dụng và quy ước trong SGK KNTT
\item (Đúng) Khi giải dạng “Tính độ biến thiên nội năng trong quá trình nhiệt động lực học”, cần xác định đúng trạng thái đầu, trạng thái cuối và quy ước dấu hoặc điều kiện áp dụng.
\item (Sai) Có thể bỏ qua mọi dữ kiện về khối lượng, nhiệt độ và hiệu suất trong mọi bài toán thuộc dạng này.
\item (Đúng) Nếu hệ nhận công 200 $J$ và tỏa 50 $J$ thì ΔU= $150\,\text{J}$.
\end{itemchoice}
}
\end{ex}

% ID: L12C1B2-19-DS
% Mức: VD
% ID: L12C1B2-19-DS
% ID: L12C1B2-19-DS
% Mức: VD
% ID: L12C1B3-103-DS
% ID: L12C1B3-103-DS
% Mức: VD
% ID: L12C1B2-85-DS
%=== Câu 47 ===%
\begin{ex}
% ID: L12C1B2-85-DS
% Mức: VD
Xét các phát biểu sau về dạng “Tính độ biến thiên nội năng trong quá trình nhiệt động lực học”
\choiceTF
{\True Khi giải dạng “Tính độ biến thiên nội năng trong quá trình nhiệt động lực học”, cần xác định đúng trạng thái đầu, trạng thái cuối và quy ước dấu hoặc điều kiện áp dụng.}
{Có thể bỏ qua mọi dữ kiện về khối lượng, nhiệt độ và hiệu suất trong mọi bài toán thuộc dạng này.}
{\True Nếu hệ nhận công 200 $J$ và tỏa 50 $J$ thì ΔU= $150\,\text{J}$.}
{Nếu A=100 $J$ và Q=-40 $J$ thì ΔU=- $140\,\text{J}$.}
\loigiai{
\begin{itemchoice}
\item (Đúng) Khi giải dạng “Tính độ biến thiên nội năng trong quá trình nhiệt động lực học”, cần xác định đúng trạng thái đầu, trạng thái cuối và quy ước dấu hoặc điều kiện áp dụng. (Vì): Khi giải dạng “Tính độ biến thiên nội năng trong quá trình nhiệt động lực học”, cần xác định đúng trạng thái đầu, trạng thái cuối và quy ước dấu hoặc điều kiện áp dụng. b) S: Có thể bỏ qua mọi dữ kiện về khối lượng, nhiệt độ và hiệu suất trong mọi bài toán thuộc dạng này. c) Đ: Nếu hệ nhận công 200 $J$ và tỏa 50 $J$ thì ΔU= $150\,\text{J}$. d) S: Nếu A=100 $J$ và Q=-40 $J$ thì ΔU=- $140\,\text{J}$. Kết luận dựa trên định nghĩa, điều kiện áp dụng và quy ước trong SGK KNTT
\item (Sai) Có thể bỏ qua mọi dữ kiện về khối lượng, nhiệt độ và hiệu suất trong mọi bài toán thuộc dạng này.
\item (Đúng) Nếu hệ nhận công 200 $J$ và tỏa 50 $J$ thì ΔU= $150\,\text{J}$.
\item (Sai) Nếu A=100 $J$ và Q=-40 $J$ thì ΔU=- $140\,\text{J}$.
\end{itemchoice}
}
\end{ex}

% ID: L12C1B2-19-TL
% Mức: NB
% ID: L12C1B2-19-TL
% ID: L12C1B2-19-TL
% Mức: NB
% ID: L12C1B3-104-TL
% ID: L12C1B3-104-TL
% Mức: NB
% ID: L12C1B2-87-TL
%=== Câu 48 ===%
\begin{bt}
% ID: L12C1B2-87-TL
% Mức: NB
Trình bày cách nhận biết và phương pháp giải dạng “Tính độ biến thiên nội năng trong quá trình nhiệt động lực học”. Phân tích nhận định sau và sửa lại nếu sai: “Nếu A=100 $J$ và Q=-40 $J$ thì ΔU=- $140\,\text{J}$.”
\loigiai{
Trước hết xác định hiện tượng, trạng thái và các đại lượng đã cho. Nêu định nghĩa hoặc hệ thức phù hợp, đổi về đơn vị SI, áp dụng đúng điều kiện và kiểm tra ý nghĩa kết quả. Nhận định cần đối chiếu: Nếu hệ nhận công 200 $J$ và tỏa 50 $J$ thì ΔU= $150\,\text{J}$. Vì vậy nhận định đã cho sai; phát biểu đúng là: Nếu hệ nhận công 200 $J$ và tỏa 50 $J$ thì ΔU= $150\,\text{J}$.
}
\end{bt}

% ID: L12C1B2-20-DS
% Mức: VDC
% ID: L12C1B2-20-DS
% ID: L12C1B2-20-DS
% Mức: VDC
% ID: L12C1B3-105-DS
% ID: L12C1B3-105-DS
% Mức: VDC
% ID: L12C1B2-88-DS
%=== Câu 50 ===%
\begin{ex}
% ID: L12C1B2-88-DS
% Mức: VDC
Xét các phát biểu sau về dạng “Tính độ biến thiên nội năng trong quá trình nhiệt động lực học”
\choiceTF
{Có thể bỏ qua mọi dữ kiện về khối lượng, nhiệt độ và hiệu suất trong mọi bài toán thuộc dạng này.}
{\True Nếu hệ nhận công 200 $J$ và tỏa 50 $J$ thì ΔU= $150\,\text{J}$.}
{Nếu A=100 $J$ và Q=-40 $J$ thì ΔU=- $140\,\text{J}$.}
{\True Khi giải dạng “Tính độ biến thiên nội năng trong quá trình nhiệt động lực học”, cần xác định đúng trạng thái đầu, trạng thái cuối và quy ước dấu hoặc điều kiện áp dụng.}
\loigiai{
\begin{itemchoice}
\item (Sai) Có thể bỏ qua mọi dữ kiện về khối lượng, nhiệt độ và hiệu suất trong mọi bài toán thuộc dạng này. (Vì): Có thể bỏ qua mọi dữ kiện về khối lượng, nhiệt độ và hiệu suất trong mọi bài toán thuộc dạng này. b) Đ: Nếu hệ nhận công 200 $J$ và tỏa 50 $J$ thì ΔU= $150\,\text{J}$. c) S: Nếu A=100 $J$ và Q=-40 $J$ thì ΔU=- $140\,\text{J}$. d) Đ: Khi giải dạng “Tính độ biến thiên nội năng trong quá trình nhiệt động lực học”, cần xác định đúng trạng thái đầu, trạng thái cuối và quy ước dấu hoặc điều kiện áp dụng. Kết luận dựa trên định nghĩa, điều kiện áp dụng và quy ước trong SGK KNTT
\item (Đúng) Nếu hệ nhận công 200 $J$ và tỏa 50 $J$ thì ΔU= $150\,\text{J}$.
\item (Sai) Nếu A=100 $J$ và Q=-40 $J$ thì ΔU=- $140\,\text{J}$.
\item (Đúng) Khi giải dạng “Tính độ biến thiên nội năng trong quá trình nhiệt động lực học”, cần xác định đúng trạng thái đầu, trạng thái cuối và quy ước dấu hoặc điều kiện áp dụng.
\end{itemchoice}
}
\end{ex}

% ID: L12C1B2-20-TLN
% Mức: NB
% ID: L12C1B2-20-TLN
% ID: L12C1B2-20-TLN
% Mức: NB
% ID: L12C1B3-106-TLN
% ID: L12C1B3-106-TLN
% Mức: NB
% ID: L12C1B2-89-TLN
%=== Câu 51 ===%
\begin{ex}
% ID: L12C1B2-89-TLN
% Mức: NB
Cho bốn nhận định: (1) Nếu hệ nhận công 200 $J$ và tỏa 50 $J$ thì ΔU= $150\,\text{J}$. (2) Nếu A=100 $J$ và Q=-40 $J$ thì ΔU=- $140\,\text{J}$. (3) Cần kiểm tra đơn vị và điều kiện áp dụng khi xử lí dạng “Tính độ biến thiên nội năng trong quá trình nhiệt động lực học”. (4) Có thể thay tùy ý nhiệt độ Celsius cho nhiệt độ tuyệt đối trong mọi công thức. Có bao nhiêu nhận định đúng? Trả lời bằng một số nguyên.
\shortans{2}
\loigiai{
Nhận định (1) và (3) đúng; (2) và (4) sai. Vậy có 2 nhận định đúng.
}
\end{ex}

% ID: L12C1B2-21-TL
% Mức: TH
% ID: L12C1B2-21-TL
% ID: L12C1B2-21-TL
% Mức: TH
% ID: L12C1B3-107-TL
% ID: L12C1B3-107-TL
% Mức: TH
% ID: L12C1B2-90-TL
%=== Câu 54 ===%
\begin{bt}
% ID: L12C1B2-90-TL
% Mức: TH
Trình bày cách nhận biết và phương pháp giải dạng “Tính độ biến thiên nội năng trong quá trình nhiệt động lực học”. Phân tích nhận định sau và sửa lại nếu sai: “Nếu hệ nhận công 200 $J$ và tỏa 50 $J$ thì ΔU= $150\,\text{J}$.”
\loigiai{
Trước hết xác định hiện tượng, trạng thái và các đại lượng đã cho. Nêu định nghĩa hoặc hệ thức phù hợp, đổi về đơn vị SI, áp dụng đúng điều kiện và kiểm tra ý nghĩa kết quả. Nhận định cần đối chiếu: Nếu hệ nhận công 200 $J$ và tỏa 50 $J$ thì ΔU= $150\,\text{J}$. Vì vậy nhận định đã cho đúng.
}
\end{bt}

% ID: L12C1B2-31-TN
% Mức: NB
% ID: L12C1B2-31-TN
% ID: L12C1B2-31-TN
% Mức: NB
% ID: L12C1B3-108-TN
% ID: L12C1B3-108-TN
% Mức: NB
% ID: L12C1B2-91-TN
%=== Câu 79 ===%
\begin{ex}
% ID: L12C1B2-91-TN
% Mức: NB
Câu 1 thuộc dạng “Tính độ biến thiên nội năng trong quá trình nhiệt động lực học”. Phát biểu nào sau đây đúng?
\choice
{Nếu A=100 $J$ và Q=-40 $J$ thì ΔU=- $140\,\text{J}$.}
{Nhiệt lượng và nhiệt độ là hai đại lượng hoàn toàn giống nhau.}
{Mọi quá trình nhiệt học đều không phụ thuộc điều kiện thực hiện.}
{\True Nếu hệ nhận công 200 $J$ và tỏa 50 $J$ thì ΔU= $150\,\text{J}$.}
\loigiai{
Phát biểu đúng là: Nếu hệ nhận công 200 $J$ và tỏa 50 $J$ thì ΔU= $150\,\text{J}$. Các phương án còn lại trái với mô hình, định nghĩa hoặc hệ thức nhiệt học tương ứng.
}
\end{ex}

% ID: L12C1B2-32-TN
% Mức: NB
% ID: L12C1B2-32-TN
% ID: L12C1B2-32-TN
% Mức: NB
% ID: L12C1B3-109-TN
% ID: L12C1B3-109-TN
% Mức: NB
% ID: L12C1B2-92-TN
%=== Câu 81 ===%
\begin{ex}
% ID: L12C1B2-92-TN
% Mức: NB
Câu 5 thuộc dạng “Tính độ biến thiên nội năng trong quá trình nhiệt động lực học”. Phát biểu nào sau đây đúng?
\choice
{Nếu A=100 $J$ và Q=-40 $J$ thì ΔU=- $140\,\text{J}$.}
{Nhiệt lượng và nhiệt độ là hai đại lượng hoàn toàn giống nhau.}
{Mọi quá trình nhiệt học đều không phụ thuộc điều kiện thực hiện.}
{\True Nếu hệ nhận công 200 $J$ và tỏa 50 $J$ thì ΔU= $150\,\text{J}$.}
\loigiai{
Phát biểu đúng là: Nếu hệ nhận công 200 $J$ và tỏa 50 $J$ thì ΔU= $150\,\text{J}$. Các phương án còn lại trái với mô hình, định nghĩa hoặc hệ thức nhiệt học tương ứng.
}
\end{ex}

% ID: L12C1B3-22-DS
% Mức: VD
% ID: L12C1B2-75-TLN
% ID: L12C1B2-75-TLN
% Mức: VD
% ID: L12C1B3-110-TLN
% ID: L12C1B3-110-TLN
% Mức: VD
% ID: L12C1B2-93-TLN
%=== Câu 83 ===%
\begin{ex}
% ID: L12C1B2-93-TLN
% Mức: VD
Một xilanh thẳng đứng chứa khí, pittông có thể trượt không ma sát. Người ta truyền nhiệt lượng $Q = 3,6\ kJ$ cho khối khí, khí giãn nở đẩy pittông đi lên. Biết áp suất khí không đổi và bằng $1,2 \times 10^5\ Pa$, thể tích khí tăng thêm $\Delta V = 8\ L$. Độ biến thiên nội năng của khối khí là bao nhiêu jun?
\shortans{2640 J}
\loigiai{
Áp dụng định luật $I$ nhiệt động lực học: $\Delta U = Q - A$, trong đó công khí thực hiện $A = p \cdot \Delta V$. Đổi đơn vị: $\Delta V = 8\ L = 8 \times 10^{-3}\ m^3$. Tính công: $A = 1,2 \times 10^5 \times 8 \times 10^{-3} = 960\ J$. Nhiệt lượng nhận vào: $Q = 3,6\ kJ = 3600\ J$ (dấu dương vì khí nhận nhiệt). Độ biến thiên nội năng: $\Delta U = 3600 - 960 = 2640\ J$. Vậy nội năng khối khí tăng $2640\ J$.
}
\end{ex}

% Mức: VD
% ID: L12C1B2-83-TN
% ID: L12C1B2-83-TN
% Mức: VD
% ID: L12C1B3-111-TN
% ID: L12C1B3-111-TN
% Mức: VD
% ID: L12C1B2-94-TN
%=== Câu 91 ===%
\begin{ex}
% ID: L12C1B2-94-TN
% Mức: VD
	Một viên đạn đại bác có khối lượng $10\ kg$ khi rơi tới đích có vận tốc $54\ km/h$. Nếu toàn bộ động năng của nó biến thành nội năng thì nhiệt lượng tỏa ra lúc va chạm vào khoảng
	\choice 
	{$1\,4580\ J$}
	{$2\,250\ J$}
	{\True $1\,125\ J$}
	{$7290\ J$}
	\loigiai{
	Theo đề bài, ta có: $Q = W_\text{đ} = \dfrac{1}{2}mv^2 = \dfrac{1}{2}10.15^2 =1 \, 125\ J$.
	}
\end{ex}

% Mức: VD
% ID: L12C1B2-84-TN
% ID: L12C1B2-84-TN
% Mức: VD
% ID: L12C1B3-112-TN
% ID: L12C1B3-112-TN
% Mức: VD
% ID: L12C1B2-95-TN
%=== Câu 92 ===%
\begin{ex}
% ID: L12C1B2-95-TN
% Mức: VD
	 Khi đun nóng một khối khí, khí giãn nở làm thể tích tăng thêm $7$ lít. Biết trong quá trình này nội năng của khí giảm $1,1\ kJ$ nhưng áp suất không đổi và bằng $3.10^5\ Pa$. Nhiệt lượng mà khối khí nhận được trong quá trình này là
	 \choice 
	 {$3,2\ kJ$}
	 {\True $1,0\ kJ$}
	 {$1,5\ kJ$}
	 {$2,7\ kJ$}
	\loigiai{
		\begin{itemize}
			\item  $\left| A\right| = p.\Delta V = 3.10^5.7.10^-3 = 2\ 100\ J \Rightarrow A = -2,1\ kJ$.
			\item $\Delta U = Q + A \Rightarrow Q = \Delta U - A = -1,1 - (-2,1) = 1,0\ kJ$.
		\end{itemize}
	}
\end{ex}

% Mức: VD
% ID: L12C1B2-85-TN
% ID: L12C1B2-85-TN
% Mức: VD
% ID: L12C1B3-113-TN
% ID: L12C1B3-113-TN
% Mức: VD
% ID: L12C1B2-96-TN
%=== Câu 93 ===%
\begin{ex}
% ID: L12C1B2-96-TN
% Mức: VD
		 Khi nung nóng một lượng không khí chứa trong một xilanh, khối khí nhận một nhiệt lượng $1,75 \ kJ$ làm nội năng của khí tăng thêm $720 \ J$. Khí giãn nở và sinh công làm pit–tông dịch chuyển. Khối khí đã thực hiện một công là
		 \choice 
		 {$2,56\ kJ$}
		 {\True $1,03\ kJ$}
		 {$1,25\ kJ$}
		 {$2,47\ kJ$}
		 \loigiai{
		 Ta có: $\Delta U = Q + A \Rightarrow A = \Delta U - Q = 0,72 - 1,75 = -1,03\ kJ \Rightarrow \left| A\right| = 1,03\ kJ$.	
		 }
		\end{ex}

% Mức: TH
% ID: L12C1B2-87-TLN
% ID: L12C1B2-87-TLN
% Mức: TH
% ID: L12C1B3-114-TLN
% ID: L12C1B3-114-TLN
% Mức: TH
% ID: L12C1B2-97-TLN
%=== Câu 95 ===%
\begin{ex}
% ID: L12C1B2-97-TLN
% Mức: TH
Một lượng khí nhận một nhiệt lượng $35,40\ kJ$ do được đun nóng, khí giãn ra và thực hiện một công $20,00\ kJ$ ra môi trường xung quanh. Nội năng của khối khí này đã biến thiên một lượng bao nhiêu kilôjun $(kJ)$ (làm tròn kết quả đến chữ số hàng phần mười)?
	\shortans[oly]{$15,4$}
	\loigiai{
	Ta có: $\Delta U = Q + A = 35,4 + (-20) = 15,4\ kJ$.
	}
\end{ex}

% Mức: VD
% ID: L12C1B2-88-TLN
% ID: L12C1B2-88-TLN
% Mức: VD
% ID: L12C1B3-115-TLN
% ID: L12C1B3-115-TLN
% Mức: VD
% ID: L12C1B2-98-TLN
%=== Câu 96 ===%
\begin{ex}
% ID: L12C1B2-98-TLN
% Mức: VD
	 Một khối khí được cung cấp nhiệt lượng $4,98\ kJ$, khí giãn nở làm tăng thể tích một lượng $\Delta V$. Trong quá trình này, nội năng của khối khí tăng $1,23\ kJ$ nhưng áp suất của khối khí không đổi và bằng $2,5.10^5\ Pa$. Giá trị của $\Delta V$ là bao nhiêu $(dm^3)$ (làm tròn kết quả đến chữ số hàng đơn vị)?
 	\shortans[oly]{$0,38$}%Thay đáp án chuẩn 
	\loigiai{
		\begin{itemize}
			\item Theo đề bài, ta có: $\Delta U = Q + A \Rightarrow A = \Delta U - Q = 1,23 - 4,98 = -3,75\ kJ$.
			\item $ \left| A\right| =p.\Delta V \Rightarrow \Delta V = \dfrac{\left| A\right|}{p} = \dfrac{3\ 750}{2,5.10^5} = 0,015\ m^3 = 15\ dm^3$.
		\end{itemize}
	
	}
\end{ex}

% Mức: VDC
% ID: L12C1B2-89-TLN
% ID: L12C1B2-89-TLN
% Mức: VDC
% ID: L12C1B3-116-TLN
% ID: L12C1B3-116-TLN
% Mức: VDC
% ID: L12C1B2-99-TLN
%=== Câu 97 ===%
\begin{ex}
% ID: L12C1B2-99-TLN
% Mức: VDC
Một vật có khối lượng $2 \, kg$ trượt không vận tốc đầu từ đỉnh xuống chân một mặt nghiêng dài $40 \, m$, nghiêng một góc $60^\circ$ so với phương ngang. Tốc độ của vật ở chân mặt phẳng nghiêng là $4,5 \, m/s$. Cho rằng, $75\%$ công của lực ma sát giữa mặt phẳng nghiêng và vật chuyển thành nội năng của vật, bỏ qua phần nhiệt lượng mặt phẳng nghiêng truyền cho vật. Lấy $g = 9,8 \, m/s^2$. Độ biến thiên nội năng của vật trong quá trình trên là bao nhiêu kilôjun (kJ) (làm tròn kết quả đến chữ số hàng phần trăm)?

	\shortans[oly]{$0,49$}%Thay đáp án chuẩn 
	\loigiai{
		\begin{itemize}
			\item Chọn mốc thế năng tại chân mặt phẳng nghiêng.
			\item Công mà vật nhận là:
			\begin{align*}
				A &= -A_{ms} = -\left(W_2 - W_1\right) = W_1 - W_2\\ 
				& = mgh-\dfrac{1}{2}mv^2 = mg\ell \sin \alpha - \dfrac{1}{2}mv^2\\
				& = 2.9,8.40.\sin 60^\circ - \dfrac{1}{2}.2.4,5^2 \approx 658,7\ J
			\end{align*}
			\item Độ biến thiên nội năng của vật: $\Delta U = 0,75.A = 0,75.658,7 \approx 494\ J = 0,494\ kJ$.
		\end{itemize}
	
		
		}
\end{ex}

% Mức: VD
% ID: L12C1B2-90-TLN
% ID: L12C1B2-90-TLN
% Mức: VD
% ID: L12C1B3-117-TLN
% ID: L12C1B3-117-TLN
% Mức: VD
% ID: L12C1B2-100-TLN
%=== Câu 98 ===%
\begin{ex}
% ID: L12C1B2-100-TLN
% Mức: VD
	\immini[thm]
	{Một xilanh có pit-tông nằm ngang như hình vẽ, xilanh chứa $500 \, cm^3$ không khí. Khi đốt nóng khí trong xilanh, khí giãn nở đẩy pit-tông qua phải làm thể tích khối khí tăng lên $720 \, cm^3$, nội năng khối khí tăng thêm $1,5 \, kJ$. Cho rằng áp suất của khối khí luôn bằng $5,8.10^5 \, Pa$. Nhiệt lượng đã cung cấp cho khối khí bao nhiêu kilôjun $(kJ)$ (làm tròn kết quả đến chữ số hàng phần trăm)?
	\shortans[oly]{$1,63$}}
	{\begin{tikzpicture}[>=stealth, thick]
			\draw [pattern=dots] (-2,1) rectangle (1,3);
			\draw [pattern=north east lines]
			(1,3) rectangle (1.2,1);
			\draw
			(2,3) -- (1,3)
			(2,1) -- (1,1)
			node [fill = white] at (-0.5,2) {Chất khí}
			node [below] at (1,1) {pit-tông};
	\end{tikzpicture}}
	
	\loigiai{
		\begin{itemize}
			\item Công do khối khí thực hiện:
			\begin{align*}
				&\left| A \right| = p.\Delta V = 5,8.10^5.\left( 720 - 500\right).10^{-6} =  127,6\ J\\
				\Rightarrow &A = - 0,1276\ kJ
			\end{align*}
			\item $\Delta U = Q + A \Rightarrow Q = \Delta U - A= 1,5 - (-0,1276) = 1,6276\ kJ \approx 1,63\ kJ$
		\end{itemize}
	}
\end{ex}
